% Ad Atticum 14.3 (ad-atticum-14-03)
% Source: https://raw.githubusercontent.com/PerseusDL/canonical-latinLit/master/data/phi0474/phi057/phi0474.phi057.perseus-lat1.xml
% Fetched by scripts/fetch_latin.py

% Dateline (Perseus): Scr. in Tusculano v Id. Apr. a. 710 (44).

\ciceroLetterOpener{CICERO ATTICO salutem}

\ciceroSection{1} Tranquillae tuae quidem litterae. quod utinam diutius! nam Matius posse negabat. ecce autem structores nostri ad frumentum profecti, cum inanes redissent, rumorem adferunt magnum Romae domum ad Antonium frumentum omne portari. πανικὸν certe; scripsisses enim. Corumbus Balbi nullus adhuc: est mihi notum nomen; bellus enim esse dicitur architectus.

\ciceroSection{2} ad obsignandum tu adhibitus non sine causa videris. volunt enim nos ita putare; nescio cur non animo quoque sentiant. sed quid haec ad nos? odorare tamen Antoni διάθεσιν ; quem quidem ego epularum magis arbitror rationem habere quam quicquam mali cogitare. tu si quid πραγματικὸν habes rescribe; sin minus, populi ἐπισημασίαν et mimorum dicta perscribito. Piliae et Atticae salutem.
